\documentclass[12pt,a4paper]{article}
\usepackage[utf8]{inputenc}
\usepackage[T1]{fontenc}
\usepackage{amsmath,amssymb}
\usepackage{graphicx,booktabs,hyperref,natbib}
\usepackage[margin=2.5cm]{geometry}

\begin{document}

\title{\bf A testable prediction for dark energy from information theory}

\author{Anonymous}
\date{June 2026}

\maketitle

\begin{abstract}
The physical nature of dark energy is unknown. The DESI collaboration has reported evidence for dynamical dark energy, but all existing explanations---quintessence, emergent dark energy, and phenomenological parameterizations---treat the equation of state $w(z)$ as a function to be fitted rather than derived. Here we show that if dark energy originates from a finite information capacity of spacetime---specifically, from the occupation of vacuum quantum information modes by cosmic structure---then $w(z)$ acquires a specific, falsifiable shape: it rises from $z=0$ to a peak at $z\sim1$--$2$, then falls back toward $w=-1$ at higher redshift. The shape is $w(z)+1\propto{\rm SFR}(z)/[H(z)\,\rho_*(z)^{n/(n+1)}]$, determined by the cosmic star formation history with a single physical parameter $n$. This prediction is non-monotonic and cannot be produced by the standard CPL parameterization or by quintessence with simple potentials. We show that the prediction is in tension with the CPL parameters preferred by DESI DR2---as expected when fitting a peaked function to a linear form---and we identify the direct measurement of $w(z)$ in redshift bins by DESI DR3 and Euclid as the decisive test. If $w(z)$ is monotonic, the model is falsified.
\end{abstract}

\section*{Introduction}

Dark energy is the largest contributor to the universe's energy budget, yet its physical nature remains unknown a quarter-century after the discovery of cosmic acceleration~\cite{riess1998,spergel2003}. The standard $\Lambda$CDM model identifies it with a cosmological constant ($w\equiv-1$). Dynamical models---quintessence~\cite{tsujikawa2013}, emergent dark energy~\cite{li2024}, and parameterized equations of state~\cite{chevallier2001,linder2003}---allow $w$ to vary with redshift but treat its functional form as an empirical choice. None provides a microphysical mechanism.

The DESI collaboration's Data Release 2 has sharpened the question~\cite{desi2025}. Baryon acoustic oscillation measurements, combined with CMB and supernova data, disfavor $\Lambda$CDM at $2.8$--$4.2\sigma$ when dynamical dark energy is allowed. In the CPL parameterization $w(a)=w_0+w_a(1-a)$~\cite{chevallier2001}, the best-fit values are $(w_0,w_a)=(-0.785\pm0.047,-0.43\pm0.10)$ for DESI+CMB+DESY5~\cite{zhang2026}. But CPL is a fitting form, not a physical theory: it describes {\it that} $w$ evolves, not {\it why}.

In this paper we take a different approach. We ask: if dark energy has an information-theoretic origin, what does that imply for $w(z)$? The answer, we find, is a specific, non-monotonic shape that is testable with near-future data---and that differs qualitatively from the CPL form.

\section*{Physical mechanism}

Consider the hypothesis that the universe possesses a finite information capacity. This idea has appeared in various guises---holographic entropy bounds~\cite{thooft1993,susskind1995}, causal set theory~\cite{sorkin2005}, and entropic gravity~\cite{verlinde2011,jacobson1995}. We adopt a specific formulation, the Discrete Graph Framework, developed in a series of works summarized in the Supplementary Information. The framework rests on two postulates:

\vspace{4pt}
\noindent{\bf P1.} Causal relations are asymmetric and directional.\\
{\bf P2.} Each minimal causal unit carries at most one bit of distinguishable information.

\vspace{4pt}
From these, the state of the universe is characterized by a scalar field $q(t)\in[0,1]$---the fraction of causal-graph nodes in the unoccupied state---whose evolution on cosmological scales is governed by a nonlinear telegraph equation (Supplementary \S2). In the slow-roll regime, and on homogeneous scales, the equation reduces to

\begin{equation}\label{eq:ode}
\gamma_0\Delta^n\frac{d\Delta}{dt} + \kappa\int_0^t\Delta\,dt' = \eta\,\dot{\rho}_*(t),
\end{equation}

where $\Delta\equiv1-q$ is the information deficit, $\gamma_0$ is a microphysical damping rate, $n$ is a damping nonlinearity index, $\kappa$ is a memory kernel strength, $\dot{\rho}_*(t)={\rm SFR}(t)$ is the cosmic star formation rate density, and $\eta$ is a coupling constant with dimensions $[M]^{-1}[L]^{3}$.

The physical picture is the following. The vacuum contains a finite number of quantum information modes---the distinct ways in which causal structure can be configured at the Planck scale. Structure formation (stars, galaxies, clusters) ``occupies'' some of these modes by establishing definite causal relations among their constituents. This occupation is irreversible on cosmological timescales: once a mode is used to encode the causal history of a structure, it is no longer available for vacuum fluctuations. The dark energy density is the energy of the remaining free modes:

\begin{equation}
\rho_{\rm DE}(t) = \rho_\Lambda\,q(t) = \rho_\Lambda\,(1-\Delta(t)).
\end{equation}

The coupling $\eta$ parameterizes how efficiently star formation occupies information modes. A microscopic calculation of $\eta$ is not available; it is calibrated from the observed $w_0$. This is the principal phenomenological input of the model.

\section*{Predicted shape of $w(z)$}

For redshifts where the memory integral in (\ref{eq:ode}) is subdominant to the star formation driving term ($z\gtrsim0.5$), the equation simplifies to $\gamma_0\Delta^n d\Delta/dt\simeq\eta\,{\rm SFR}$. Integrating yields $\Delta^{n+1}\propto\rho_*(t)$, where $\rho_*(t)=\int_0^t{\rm SFR}(t')dt'$ is the cumulative stellar mass density. From energy conservation, the equation of state is

\begin{equation}\label{eq:shape}
\boxed{w(z)+1 = C\cdot\frac{{\rm SFR}(z)}{H(z)\,\rho_*(z)^{\,n/(n+1)}}\cdot\frac{1}{1-\Delta(z)}},
\end{equation}

where $C=\frac{1}{3}\sqrt{\eta/(2\gamma_0)}$ is calibrated from $w_0$. The factor $(1-\Delta)^{-1}$ is of order unity for the calibrated parameters. The normalized shape

\begin{equation}\label{eq:norm}
\boxed{S(z)\equiv\frac{w(z)+1}{w(0)+1} = \frac{{\rm SFR}(z)/[H(z)\rho_*(z)^{\,n/(n+1)}]}{{\rm SFR}(0)/[H_0\rho_{*,0}^{\,n/(n+1)}]}}
\end{equation}

is independent of the coupling constant $C$ and depends only on known astrophysical quantities and the single parameter $n$.

The key feature of equation (\ref{eq:shape}) is that $w(z)+1$ inherits the shape of the cosmic star formation history. The star formation rate density rises from $z=0$ to a broad peak at $z\sim1$--$2$~\cite{madau2014}, then declines at higher redshift. Consequently, $w(z)+1$ is predicted to have a peak. For the natural choice $n=1$, which corresponds to damping proportional to the information deficit (the simplest mean-field expectation, analogous to the Drude model for resistivity), the peak is at $z\approx1.2$ with $w_{\rm peak}\approx-0.3$ (Fig.~\ref{fig:wz}a). The peak location shifts with $n$: larger $n$ moves the peak to lower redshift and reduces its amplitude.

\begin{figure}[htbp]
\centering
\includegraphics[width=\textwidth]{../figures/fig2_wz_shape.png}
\caption{\bf Predicted $w(z)$ shape.} {\bf a}, $w(z)$ for $n=1.0$ (blue), $n=0.7$ (dashed green), and $n=1.3$ (dashed red). The blue band shows the theory uncertainty from the $\pm27\%$ star formation rate and $\pm11\%$ stellar mass density systematics. The red star marks the DESI $z\approx0$ constraint. $\Lambda$CDM is the dotted line. The peak---the model's central prediction---is at $z\sim1$--$2$. {\bf b}, $w(z)+1$ on a logarithmic scale. The shape at $z>2$ (dashed) is sensitive to the memory kernel form.}
\label{fig:wz}
\end{figure}

\section*{Tension with the CPL parameterization}

The CPL form $w(a)=w_0+w_a(1-a)$ is linear in $(1-a)$ and therefore monotonic in redshift. A peaked $w(z)$ cannot be represented by CPL for any choice of $(w_0,w_a)$. When a peaked function is projected onto CPL by least-squares fitting, the resulting $w_a$ is positive---reflecting the rising portion of the peak---while the DESI data prefer $w_a<0$~\cite{zhang2026}. This is a real tension, but it is a tension between a peaked prediction and a linear fitting form, not necessarily between the model and the underlying distance data.

We illustrate this in Fig.~\ref{fig:cpl_vs_dgf}. The DGF $w(z)$ (blue curve) rises from $z=0$ to a peak at $z\sim1.2$, then falls. The CPL best-fit to DESI data (red dashed) is a monotonic straight line in $(1-a)$. The two functions differ qualitatively. The CPL projection of DGF gives $w_a>0$ for all physically motivated fitting ranges (Table~\ref{tab:proj}), while DESI prefers $w_a<0$. This sign difference is the expected consequence of fitting a peaked function to a linear parameterization; it does not indicate that DGF is incompatible with the DESI distance measurements, which must be tested directly.

\begin{figure}[htbp]
\centering
\includegraphics[width=0.85\textwidth]{../figures/figS1_cpl_vs_dgf.png}
\caption{\bf DGF $w(z)$ compared to the CPL best-fit.} The DGF prediction (blue) has a peak. The CPL form preferred by DESI (red dashed) is monotonic. No choice of $(w_0,w_a)$ in the CPL parameterization can produce a peak. The functions differ qualitatively, so their CPL parameters differ---this is expected, not a contradiction.}
\label{fig:cpl_vs_dgf}
\end{figure}

\begin{table}[htbp]
\centering
\caption{\bf CPL projection of DGF $w(z)$ for different fitting ranges ($n=1$).}
\label{tab:proj}
\begin{tabular}{lccc}
\toprule
Fitting range $z\in$ & $w_0$ (CPL) & $w_a$ (CPL) & Sign of $w_a$ \\
\midrule
$[0,0.5]$  & $-0.82$ & $+0.94$ & Positive \\
$[0,1.0]$  & $-0.91$ & $+1.45$ & Positive \\
$[0,1.5]$  & $-1.03$ & $+1.90$ & Positive \\
$[0,2.0]$  & $-1.09$ & $+2.09$ & Positive \\
$[0,2.33]$ & $-1.08$ & $+2.06$ & Positive \\
$[0,3.0]$  & $-0.96$ & $+1.78$ & Positive \\
\midrule
DESI DR2 & $-0.785\pm0.047$ & $-0.43\pm0.10$ & Negative \\
\bottomrule
\end{tabular}
\end{table}

\section*{The proper test}

The CPL projection tension means that the correct comparison between DGF and DESI data cannot be performed in the $(w_0,w_a)$ plane. The model predicts a specific $w(z)$ shape; the data constrain distance-redshift relations $D_A(z)$ and $H(z)$. The proper test is to compute the DGF distance-redshift relation directly from the Friedmann equation with the DGF $w(z)$ inserted, and to compare it against the DESI BAO distance measurements.

We have performed this calculation (Supplementary \S8). For $n=1$, the DGF distance-redshift relation differs from $\Lambda$CDM by $1$--$2\%$ at $z\sim1$, which is within the sensitivity range of DESI DR2 BAO. A full likelihood analysis against the published DESI BAO data vector, marginalized over $n$ and the coupling $\eta/\gamma_0$, is beyond the scope of this paper and will be reported separately. At the current stage, what we can state is that the DGF $w(z)$ shape makes a definite prediction---the peak---and that DESI DR3 and Euclid will provide the data to test it.

\section*{Discussion}

The central result of this paper is a prediction: if dark energy originates from the finite information capacity of spacetime, then $w(z)$ has a peak at $z\sim1$--$2$. This prediction follows from the cosmic star formation history, which is measured independently of dark energy. It is non-monotonic and cannot be mimicked by the CPL parameterization or by quintessence with simple potentials. It is falsifiable by DESI DR3 and Euclid within the next several years.

We have been explicit about the limitations of the present analysis. The coupling $\eta$ between star formation and information-mode occupation is calibrated from $w_0$, not computed from first principles. The theory uncertainty is currently dominated by the $\pm27\%$ systematic error in the star formation rate normalization~\cite{madau2014}, which is comparable to the DESI measurement precision. The CPL projection of a peaked $w(z)$ necessarily yields $w_a>0$, in tension with the DESI-preferred $w_a<0$; this tension is a consequence of the inadequacy of CPL for describing non-monotonic equations of state, but it underscores that only a direct $w(z)$ or $D(z)$ comparison can adjudicate between the model and the data.

The broader DGF program aims to derive quantum mechanics, general relativity, and black hole thermodynamics from the same information-theoretic postulates. These derivations are summarized in the Supplementary Information; each is at a different stage of completion. The cosmological prediction presented here is independent of those developments and stands or falls on its own.

The peak in $w(z)$ at $z\sim1$--$2$ is the decisive test. If upcoming surveys find $w(z)$ to be monotonic---specifically, if binned equation-of-state measurements from DESI DR3 and Euclid are consistent with a linear function in $(1-a)$ at $>3\sigma$---the model presented here is falsified. If, instead, those measurements reveal a peak, they will point toward an information-theoretic origin for dark energy.

\section*{Methods}

\subsection*{Star formation history}
We use the Madau \& Dickinson (2014)~\cite{madau2014} fit: ${\rm SFR}(z)=0.015(1+z)^{2.7}/[1+((1+z)/2.9)^{5.6}]\,M_\odot\,{\rm yr}^{-1}\,{\rm Mpc}^{-3}$. The cumulative stellar mass density is $\rho_*(z)=0.72\int_z^\infty{\rm SFR}(z')|dt/dz'|dz'$, with the factor $0.72$ accounting for mass return (Chabrier IMF). The local SFR normalization uncertainty is $\pm27\%$; the $\rho_{*,0}$ uncertainty is $\pm11\%$~\cite{driver2018}. Robustness to the choice of SFR compilation is checked using Behroozi et al.\ (2019)~\cite{behroozi2019} (Supplementary Fig.~S2).

\subsection*{Cosmology}
Planck 2018~\cite{planck2018}: $H_0=67.4$ km s$^{-1}$ Mpc$^{-1}$, $\Omega_m=0.315$, $\Omega_\Lambda=0.685$.

\subsection*{Equation of state}
We solve equation (\ref{eq:ode}) in the driving-dominated regime, coupled to the Friedmann equation via Picard iteration. Convergence (max$|\delta H/H|<10^{-6}$) is achieved in 3--5 iterations. The coupling $C$ is calibrated from $w_0$; results are insensitive to the choice of $w_0$ within the DESI $1\sigma$ range.

\subsection*{CPL projection}
Weighted least-squares fit of $w(a)=w_0+w_a(1-a)$ to the DGF $w(z)$ over the specified redshift range. Weights $\propto(1+z)^2/H(z)$ approximate the DESI BAO sensitivity.

\subsection*{DESI constraints}
The DESI+CMB+DESY5 joint constraint~\cite{zhang2026} is used for the $(w_0,w_a)$ comparison: $w_0=-0.785\pm0.047$, $w_a=-0.43\pm0.10$, $\rho=0.315$. For the direct $D(z)$ comparison (Supplementary \S8), we use the DESI DR2 BAO distance measurements from~\cite{desi2025}.

\subsection*{Code and data availability}
All numerical code, figure data, and the machine-readable data package are available at the project repository (Python~3.12, NumPy, SciPy, Matplotlib).

\begin{thebibliography}{99}

\bibitem{riess1998}
Riess, A.G.\ et al.\ Observational evidence from supernovae for an accelerating universe. {\it Astron.\ J.} {\bf 116}, 1009--1038 (1998).

\bibitem{spergel2003}
Spergel, D.N.\ et al.\ First-year WMAP observations: cosmological parameters. {\it Astrophys.\ J.\ Suppl.} {\bf 148}, 175--194 (2003).

\bibitem{tsujikawa2013}
Tsujikawa, S.\ Quintessence: a review. {\it Class.\ Quant.\ Grav.} {\bf 30}, 214003 (2013).

\bibitem{li2024}
Li, X.\ et al.\ Emergent dark energy. {\it Phys.\ Rev.\ D} {\bf 109}, 043510 (2024).

\bibitem{chevallier2001}
Chevallier, M.\ \& Polarski, D.\ Accelerating universes with scaling dark matter. {\it Int.\ J.\ Mod.\ Phys.\ D} {\bf 10}, 213--223 (2001).

\bibitem{linder2003}
Linder, E.V.\ Exploring the expansion history of the universe. {\it Phys.\ Rev.\ Lett.} {\bf 90}, 091301 (2003).

\bibitem{desi2025}
DESI Collaboration.\ DESI DR2 results II: BAO measurements and cosmological constraints. arXiv:2503.14738 (2025).

\bibitem{zhang2026}
Zhang, X.\ et al.\ Robust evidence for dynamical dark energy from DESI DR2, CMB, and supernovae. SSRN 6215384 (2026).

\bibitem{thooft1993}
't Hooft, G.\ Dimensional reduction in quantum gravity. arXiv:gr-qc/9310026 (1993).

\bibitem{susskind1995}
Susskind, L.\ The world as a hologram. {\it J.\ Math.\ Phys.} {\bf 36}, 6377--6396 (1995).

\bibitem{sorkin2005}
Sorkin, R.D.\ Causal sets: discrete gravity. In {\it Lectures on Quantum Gravity} (Springer, 2005).

\bibitem{verlinde2011}
Verlinde, E.\ On the origin of gravity and the laws of Newton. {\it J.\ High Energy Phys.} {\bf 2011}, 29 (2011).

\bibitem{jacobson1995}
Jacobson, T.\ Thermodynamics of spacetime. {\it Phys.\ Rev.\ Lett.} {\bf 75}, 1260--1263 (1995).

\bibitem{madau2014}
Madau, P.\ \& Dickinson, M.\ Cosmic star-formation history. {\it Annu.\ Rev.\ Astron.\ Astrophys.} {\bf 52}, 415--486 (2014).

\bibitem{planck2018}
Planck Collaboration.\ Planck 2018 results. VI. {\it Astron.\ Astrophys.} {\bf 641}, A6 (2020).

\bibitem{driver2018}
Driver, S.P.\ et al.\ GAMA/G10-COSMOS/3D-HST. {\it Mon.\ Not.\ R.\ Astron.\ Soc.} {\bf 475}, 2891--2935 (2018).

\bibitem{behroozi2019}
Behroozi, P.\ et al.\ UniverseMachine. {\it Mon.\ Not.\ R.\ Astron.\ Soc.} {\bf 488}, 3143--3194 (2019).

\end{thebibliography}

\end{document}
